% ─────────────────────────────────────────────────────────────────────────────
This chapter presents the current state of the art across the four research areas
that underpin BookForMe: (1) agentic AI systems built on large language models;
(2) task-oriented dialogue and state tracking; (3) bilingual and code-switched NLU
for South Asian languages; and (4) distributed booking systems with concurrency
guarantees.  It also establishes the novelty of our work by identifying gaps that
existing research does not address and that directly motivate our design choices.

% ─────────────────────────────────────────────────────────────────────────────
\section{From Chatbots to Autonomous Agents}
\label{sec:lit-agents}

\subsection{Defining the Agentic Architecture}

The term ``agent'' has been used broadly in software engineering and robotics, but
its application to large language models (LLMs) represents a specific architectural
shift.  Wang et al.~\cite{wang2024} propose a unified framework in which an LLM-based
agent comprises four modules: \textit{Profiling} (role and persona definition),
\textit{Memory} (short-term context and long-term persistent state), \textit{Planning}
(decomposing goals into subgoals and choosing actions), and \textit{Action} (calling
external tools and APIs to effect changes in the environment).  In this framework the
LLM functions as a ``Brain'' that orchestrates a Perception $\rightarrow$ Reasoning
$\rightarrow$ Action cycle, fundamentally distinguishing it from a passive token
predictor.

Xi et al.~\cite{xi2023} further distinguish agents from traditional conversational
systems by their capacity to \textit{initiate and manage workflows} rather than merely
respond to prompts.  A conventional chatbot can answer ``Is the 5~PM slot available?'';
our agent must answer ``Book the 5~PM slot, lock it for ten minutes, verify the advance
payment, and notify the vendor.''  The latter requires stateful read--write operations
across multiple external systems---precisely the capability gap this work fills.

\subsection{Cognitive Architectures: ReAct, PRACT, and Reflexion}

Yao et al.~\cite{yao2023} introduced the \textbf{ReAct} (Reasoning + Acting) paradigm,
which interleaves chain-of-thought reasoning with external actions, enabling explicit
planning traces such as:
\textit{Thought}: User requests 5~PM padel.
$\rightarrow$ \textit{Action}: \texttt{get\_availability("padel","2026-03-15","17:00")}.
$\rightarrow$ \textit{Observation}: Slot available.
$\rightarrow$ \textit{Action}: \texttt{create\_booking(user, slot)}.

Liu et al.~\cite{liu2024} extend this with the \textbf{PRACT} agent, adding a
self-reflection step \textit{before} executing irreversible actions.  Before calling
\texttt{create\_booking()}, the agent verifies that the action is consistent with the
user's latest stated intent---a critical safeguard where a confirmed booking constitutes
a financial commitment.

Shinn et al.~\cite{shinn2023} demonstrate in \textbf{Reflexion} that verbal
reinforcement and persistent state across iterations significantly improve agent
reliability on multi-step tasks.  The agent receives textual feedback from the
environment (e.g., ``payment amount incorrect'') and updates its policy accordingly
without gradient updates.  This pattern directly informs our OCR rejection loop, where
the agent re-prompts the user upon detecting an incorrect payment amount.

\subsection{Orchestration Frameworks: LangChain and LangGraph}

LangChain~\cite{mavroudis2024} provides a modular abstraction layer that decouples
application logic from specific LLM providers, enabling composition of prompt
templates, vector stores, and output parsers into unified pipelines.  However,
standard LangChain implementations rely on Directed Acyclic Graphs (DAGs), which
cannot represent the iterative, self-correcting loops required for asynchronous
WhatsApp conversations where users may reply late, change their mind, or submit
incorrect payment proofs.

\textbf{LangGraph}~\cite{langgraph2024} addresses this by replacing DAGs with
\textit{Cyclic State Graphs}: the agent's behaviour is modelled as a state machine
where nodes represent actions and edges encode conditional logic.  The framework
provides loop control, state persistence across invocations, and fault recovery---
all essential for BookForMe's multi-turn WhatsApp booking conversations.
Shang et al.~\cite{shang2024} confirm in AgentSquare that modular architectures
separating Planning, Reasoning, and Memory are substantially more robust for complex
tasks than monolithic prompt designs.

\subsection{Multi-Agent Decomposition}

H{\"a}ndler~\cite{handler2023} proposes a multi-dimensional taxonomy of autonomous
LLM-powered multi-agent architectures, showing that complex workflows benefit from
assigning distinct roles to specialised sub-agents (intent parsing, verification,
database updates, external communications) coordinated by a central orchestrator.
Wu et al.~\cite{wu2023autogen} in AutoGen further argue that effective problem-solving
requires cyclic loops---the ability to retry failed actions, request missing
information, and self-correct---rather than purely linear pipelines.  BookForMe
adopts this multi-node, cyclic design in its LangGraph implementation.

% ─────────────────────────────────────────────────────────────────────────────
\section{Task-Oriented Dialogue and State Tracking}
\label{sec:lit-tod}

Task-Oriented Dialogue (TOD) systems have evolved from modular pipelines (NLU $\to$
Dialogue State Tracking $\to$ Policy $\to$ NLG) to end-to-end architectures.
\textbf{SimpleTOD}~\cite{hosseini2020} treats the entire dialogue as unified causal
language modelling, jointly predicting belief states, system actions, and responses
in a single autoregressive model.  While effective on standard benchmarks (MultiWOZ),
it requires large annotated corpora---a significant barrier for our domain where no
pre-built dataset exists.

Shin et al.~\cite{shin2022} address this data scarcity problem with \textbf{DS2}
(Dialogue Summaries as Dialogue States), converting conversation history into
templated natural-language summaries that map directly to structured slot
representations.  This summarisation-based approach is particularly useful for
short, noisy WhatsApp messages where maintaining a full belief-state vector is
impractical.

% ─────────────────────────────────────────────────────────────────────────────
\section{Bilingual and Code-Switched NLU}
\label{sec:lit-bilingual}

\subsection{Roman Urdu Challenges}

Processing Roman Urdu introduces challenges largely absent from standard multilingual
NLP.  Bhat et al.~\cite{bhat2023} classify Roman Urdu as exhibiting severe
\textit{orthographic variation}: a single word can be rendered in multiple phonetically
plausible romanisations by different writers (e.g., \textit{waqt}, \textit{vakt},
\textit{wkht}, \textit{waqat}, and \textit{tym} all mean ``time'').  Standard subword
tokenisers (BPE, WordPiece) trained on formal corpora frequently fail to map these
variants to a shared semantic representation.  Integrating a normalisation layer
before intent classification is therefore necessary.

\subsection{Cross-Lingual Transfer and Adapters}

Yu et al.~\cite{yu2023} address cross-lingual dialogue state tracking via contrastive
learning, aligning slot embeddings across languages so that semantically equivalent
slots in different languages map to nearby vector representations.  Wu et al.~\cite{wu2023adapter}
show that parameter-efficient adapter methods can transfer pretrained models to
code-switched contexts without full retraining, significantly reducing data and
compute requirements.

The Qwen2.5 technical report~\cite{qwen2025} demonstrates that the 7B variant
achieves 79.3 on the MMLU-Multi-Understanding benchmark, the highest among models
of comparable size, validating its use as a multilingual backbone.

\subsection{Multilingual Pretrained Models}

Devlin et al.~\cite{devlin2019} introduced \textbf{BERT}, establishing bidirectional
masked language modelling as the dominant pretraining paradigm.  The multilingual
variant (mBERT) was pretrained on 104 languages, including Urdu, providing limited
but useful cross-lingual transfer.  Conneau et al.~\cite{conneau2020} demonstrate
that \textbf{XLM-RoBERTa}, trained on a substantially larger multilingual corpus
with improved data filtering, outperforms mBERT on cross-lingual classification.
Sanh et al.~\cite{sanh2019} show that \textbf{DistilBERT}, a distilled 66M-parameter
version of BERT, retains 97\% of BERT's performance while being 40\% smaller and
60\% faster---a highly attractive tradeoff for latency-sensitive production deployment.

\subsection{Data Augmentation for Low-Resource Settings}

Given the absence of pre-built bilingual booking datasets, data augmentation is
essential.  Hou et al.~\cite{hou2020} show that back-translation and contextual
augmentation improve slot-filling F1 by 6--17\% and intent accuracy by up to 12\%
in booking domains with fewer than 500 annotated utterances.  An et al.~\cite{an2022}
demonstrate that controllable T5-based generation yields 8--18\% absolute gains in
intent classification under 10-shot conditions for restaurant and hotel reservation
tasks.  He et al.~\cite{he2024} in \textbf{SynthDST} show that fully synthetic
LLM-generated dialogues can surpass real-data baselines, achieving 4--12\% higher
Joint Goal Accuracy on MultiWOZ when human annotations are extremely limited.
BookForMe leverages LLM-driven paraphrasing and synthetic generation following
these findings.

% ─────────────────────────────────────────────────────────────────────────────
\section{Retrieval, Planning, and Tool Use}
\label{sec:lit-tools}

Schick et al.~\cite{schick2023} demonstrated in \textbf{Toolformer} that LLMs can
learn to invoke external APIs through self-supervised training, deciding autonomously
when and how to call tools---a key capability for our agent's database interactions.
Patil et al.~\cite{patil2023} (\textbf{Gorilla}) and Qin et al.~\cite{qin2023}
(\textbf{ToolLLM}) extend this to controlled invocation of structured APIs, providing
the theoretical basis for our \texttt{get\_availability()},
\texttt{create\_booking()}, and \texttt{get\_services()} tool set.

Lewis et al.~\cite{lewis2020} introduced \textbf{Retrieval-Augmented Generation
(RAG)}, conditioning generation on dynamically retrieved documents.  In BookForMe,
RAG enables the agent to access up-to-date vendor availability and pricing from
Firestore rather than relying on memorised parametric knowledge---critical for a
real-time booking system.  Chen et al.~\cite{chen2025} present adaptive retrieval
strategies where the agent decides when retrieval is necessary, minimising
unnecessary database queries and reducing latency.

% ─────────────────────────────────────────────────────────────────────────────
\section{Distributed Booking and Concurrency Safety}
\label{sec:lit-concurrency}

Mong'are~\cite{mongare2024} provides a thorough analysis of idempotency in APIs
and distributed systems, specifically addressing the double-booking problem that
arises when concurrent requests modify shared mutable state.  He recommends
idempotent APIs combined with Optimistic Concurrency Control (OCC) or distributed
locking as the preferred approach for booking systems.  Our implementation directly
follows this recommendation, using Firestore atomic transactions with a
\texttt{hold\_expires\_at} timestamp to implement OCC.

Singh et al.~\cite{singh2000} demonstrate the effectiveness of reinforcement learning
for dialogue policy optimisation in the NJFun system, working with compressed 62-state
state spaces to ``greatly reduce'' training data requirements---an approach that
validates BookForMe's data-efficient learning strategy for asynchronous multi-party
interactions.

% ─────────────────────────────────────────────────────────────────────────────
\section{Novelty and Gap Analysis}
\label{sec:lit-gap}

Table~\ref{tab:lit-gap} summarises the specific gaps that BookForMe addresses
relative to the existing literature and systems.

\begin{table}[H]
\centering
\caption{Gap analysis: BookForMe versus existing dialogue and booking paradigms.}
\label{tab:lit-gap}
\small
\begin{tabularx}{\textwidth}{>{\raggedright\arraybackslash}p{0.20\textwidth}YC{0.10\textwidth}C{0.12\textwidth}C{0.12\textwidth}C{0.12\textwidth}}
\toprule
\textbf{System} & \textbf{Focus} & \textbf{Tool Use} &
\textbf{Roman Urdu} & \textbf{Async Mem.} & \textbf{Txn.\ Safety} \\
\midrule
GPT / General LLM    & General chat        & \cmark & \cmark (partial) & \xmark & \xmark \\
SimpleTOD            & Research datasets   & \xmark & \xmark           & \xmark & \xmark \\
ReAct Agents         & Web automation      & \cmark & \xmark           & \cmark (limited) & \xmark \\
Domain-specific apps & Single category     & \xmark & \xmark           & \xmark & \cmark (partial) \\
\textbf{BookForMe}   & \textbf{WA Booking} & \cmark & \cmark           & \cmark & \cmark \\
\bottomrule
\end{tabularx}
\end{table}

\noindent The primary research gaps addressed are:
\begin{itemize}
  \item \textbf{Roman Urdu code-switching:} Existing DST and NLU research does not
        target informal WhatsApp-style Roman Urdu, presenting significant gaps in
        language modelling and orthographic normalisation.
  \item \textbf{Transactional booking guarantees:} While agent tool-use is well
        studied, very few works address concurrency, idempotency, and transactional
        safety in real-world booking systems.
  \item \textbf{Asynchronous conversation flow:} Most TOD frameworks assume
        synchronous interaction.  WhatsApp requires persistent memory and session
        resumption across arbitrary time delays.
  \item \textbf{Evaluation for informal-economy SMEs:} Agent evaluation literature
        does not address booking systems in informal economies where reliability, low
        latency, and no-infrastructure vendor onboarding are essential.
\end{itemize}

BookForMe synthesises the reviewed literature into a practical LLM-driven booking
agent that is novel precisely because it operates at the intersection of all four
gaps listed above.

%%% Local Variables:
%%% mode: latex
%%% TeX-master: "../report"
%%% End:
